\pagenumbering{roman}

\cleardoublepage
\phantomsection
\addcontentsline{toc}{chapter}{Arabic Abstract}
\begin{otherlanguage}{arabic}
{\Huge\bfseries ملخص\par}
\vspace{1.5cm}
أدى التحول الرقمي في سوق العمل إلى زيادة عدد عروض التدريب المنشورة عبر الإنترنت. غير أن هذا التوسع يطرح صعوبات في إدارة البيانات، لأن العروض تكون غالبا غير مهيكلة، وغير متجانسة، وموزعة بين منصات متعددة. تعتمد كثير من أنظمة التوظيف التقليدية على المطابقة اللفظية المباشرة، وهو ما يحد من قدرتها على معالجة الفجوة الدلالية بين مسميات الوظائف، ووصف المهارات، والصيغ المختلفة التي يستعملها أصحاب العمل. وفي المقابل، فإن الاعتماد الواسع على النماذج اللغوية الكبيرة قد يسبب تكلفة مالية وزمنا كبيرا للمعالجة عند التعامل مع تدفق مستمر وكبير من البيانات. يقترح هذا العمل خط معالجة دلاليا خفيفا يهدف إلى تحويل بيانات التدريب غير المنتظمة إلى قاعدة معرفة منظمة وقابلة للبحث. تجمع المنهجية بين الاستخراج الحتمي القائم على القواعد، والنمذجة الخفيفة باستعمال نماذج \textLR{Transformer} مثل \textLR{all-MiniLM-L6-v2 transformer}، والمواءمة مع تصنيف \textLR{ESCO} المهني، مع استعمال محدود للنماذج اللغوية في الحالات الغامضة فقط. كما يتضمن النظام آلية لإدخال ملفات \textLR{JSON} متعددة البنى اعتمادا على المطابقة التقريبية بين المفاتيح والاستدلال من المحتوى. تحفظ البيانات المعالجة في بنية بحث هجينة تستعمل فهرساً شعاعياً في الذاكرة والبحث اللفظي. يعتمد التقييم المعروض في هذا البحث على مجموعة بيانات حقيقية تضم \textLR{327} عرض تدريب. وتشير الأدلة المعروضة إلى أن البحث الهجين يمكن أن يحسن الملاءمة والاسترجاع مقارنة بالبحث المعتمد على الكلمات المفتاحية فقط، مع الحفاظ على بنية عملية قابلة للتشغيل على عتاد عادي. وبذلك تتمثل مساهمة البحث في تقديم إطار عملي لتنظيم عروض التدريب دلاليا، مع تحديد عناصر المقارنة والنشر التي تحتاج إلى استكمال لاحق.
\end{otherlanguage}

\chapter*{Abstract}
\addcontentsline{toc}{chapter}{Abstract}
The digital transformation of the labor market has increased the number of internship listings published online. However, this expansion also creates data-management difficulties because listings are frequently unstructured, inconsistent, and distributed across multiple platforms. Traditional recruitment systems often rely on lexical matching, which has limited ability to address the semantic gap between job titles, skill descriptions, and employer-specific terminology. Conversely, approaches that depend primarily on Large Language Models (LLMs) may introduce high financial cost and processing latency when applied to continuous large-scale ingestion. This thesis proposes a lightweight, tiered approach for transforming noisy internship data into a structured and searchable knowledge base. The implemented prototype combines deterministic rule-based cleaning, salary parsing, skill extraction, category assignment, stable identifier generation, JSON-based persistence, and a Next.js interface for browsing and inspecting internship records. The broader architecture also defines an extension path toward lightweight transformer mapping, ESCO alignment, and hybrid lexical-semantic retrieval for cases where keyword heuristics are insufficient. The evaluation reported in the thesis uses a real-world dataset of 327 listings. The contribution of the thesis is therefore a practical framework for internship-data organization that demonstrates a low-cost structured catalog and identifies the semantic retrieval, external comparison, and deployment evidence that should be completed in future work.

\chapter*{Résumé}
\addcontentsline{toc}{chapter}{Résumé}
La transformation numérique du marché du travail a augmenté le nombre d'offres de stage publiées en ligne. Toutefois, cette expansion crée des difficultés de gestion des données, car les offres sont souvent non structurées, hétérogènes et réparties entre plusieurs plateformes. Les systèmes de recrutement fondés principalement sur la correspondance lexicale restent limités face à l'écart sémantique entre les intitulés de postes, les descriptions de compétences et les formulations propres aux employeurs. À l'inverse, les approches reposant largement sur de grands modèles de langage peuvent entraîner un coût financier et une latence élevés lorsqu'elles sont appliquées à une ingestion continue à grande échelle. Ce mémoire propose une approche légère pour transformer des données de stages bruitées en une base structurée et interrogeable. Le prototype réalisé combine le nettoyage déterministe, l'extraction de salaires, l'extraction de compétences par lexique, la catégorisation des postes, la génération d'identifiants stables, la persistance en JSON et une interface Next.js pour consulter les offres. L'architecture générale définit aussi une évolution possible vers la normalisation par modèles Transformer, l'alignement avec ESCO et la recherche hybride lexicale-sémantique lorsque les heuristiques par mots-clés deviennent insuffisantes. L'évaluation rapportée dans le mémoire s'appuie sur un ensemble réel de 327 offres. La contribution principale du mémoire est donc un cadre pratique pour organiser les offres de stage à faible coût, tout en identifiant les éléments de recherche sémantique, de comparaison externe et de déploiement qui restent à compléter.

\chapter*{Acknowledgements}
\addcontentsline{toc}{chapter}{Acknowledgements}
We would like to express our deepest gratitude to our supervisor, Mohamed Taki Eddine SEDDIK, for his invaluable guidance, constant encouragement, and insightful feedback throughout the development of this thesis. His supervision was instrumental in shaping the direction of this research.

We are also grateful to the members of the Department of Computer Science at the University of Batna 2 for providing a supportive academic environment and for the knowledge shared during our Master’s studies.

Special thanks go to our colleagues and friends who offered support, proofread chapters, and engaged in discussions that helped refine the ideas presented in this work. Finally, we would like to thank our family for their unconditional support, patience, and belief in us. This journey would not have been possible without them.

\tableofcontents
\listoffigures

\printglossaries
\clearpage
\pagenumbering{arabic}
